% a06_BALtorus_b: part 3 of the appendix app:a06_BALtorus (input from sec/a06_BALtorus.tex).
% The fixed-spacing volume limit (lem:s10:V) and the c-number lemma (hyp:s03:T2).

\subsection{The volume limit at fixed spacing}\label{ssec:a06:V}

All quantities of \cref{lem:s10:V} are built from finitely many torus Riemann sums $\RT F$ of smooth periodic functions, and the volume limit at fixed $\eps$ is the convergence of Riemann sums to integrals. We record the rate, which is also used in \cref{ssec:a06:T2}.

\begin{lemma}[torus Riemann sums]\label{lem:a06:RS}
Let $d_0=3L_0$, let $F\in C^\infty(\R^{d_0})$ be $(2\pi/\eps)$-periodic in every coordinate, and let $\RT F$ be its Riemann sum over $(T^*)^{L_0}$ for a torus of side $\ell\in\eps\N$. For every $r\in\N$,
\[
\Big|\RT F-\int_{\BZ^{L_0}}F\Big|\ \le\ \sum_{n\in\Z^{d_0}\setminus\{0\}}\Big(\frac{\pi}{2\ell|n|}\Big)^r\,\big\|\partial^r_{n/|n|}F\big\|_{L^1(\BZ^{L_0})},
\]
$\partial_v$ the directional derivative; in particular, for $r\ge d_0+1$, $|\RT F-\int F|\le C_{d_0,r}\,\ell^{-r}\max_{|\alpha|=r}\|\partial^\alpha F\|_{L^1(\BZ^{L_0})}$ with $C_{d_0,r}=(\pi/2)^rd_0^{r/2}\sum_{n\ne0}|n|^{-r}$. If $\int_{\BZ^{L_0}}F=0$, the same bounds hold for $|\RT F|$.
\end{lemma}

\begin{proof}
Expand $F$ in its Fourier series on the period cell, $F(p)=\sum_{j}c_je^{i\eps j\cdot p}$. With $h=2\pi/\ell$ and $N_\ell=\ell/\eps$, $\sum_{p\in(T^*)^{L_0}}e^{i\eps j\cdot p}=N_\ell^{d_0}\mathbf 1[j\in N_\ell\Z^{d_0}]$, so $\RT F=\sum_{n\in\Z^{d_0}}\hat F(n)$ with $\hat F(n)\defeq\int_{\BZ^{L_0}}F(p)e^{-i\ell n\cdot p}dp$ (Poisson summation); the term $n=0$ is $\int F$. For $n\ne0$ put $\delta\defeq\pi n/(\ell|n|^2)$; periodicity gives $\int F(p-\delta)e^{-i\ell n\cdot p}dp=e^{-i\ell n\cdot\delta}\hat F(n)=-\hat F(n)$, so $2\hat F(n)=\int(F-F(\cdot-\delta))e^{-i\ell n\cdot p}$ and, iterating $r$ times, $|\hat F(n)|\le2^{-r}\|\Delta^r_\delta F\|_1$. The difference quotient is $\Delta_\delta G(p)=\int_0^1(\delta\cdot\nabla G)(p-\theta\delta)\,d\theta$, so by Fubini and the shift invariance of the $L^1$ norm on the period cell $\|\Delta_\delta G\|_1\le|\delta|\,\|\partial_{\delta/|\delta|}G\|_1$; applied $r$ times in the same direction, $|\hat F(n)|\le(\pi/(2\ell|n|))^r\|\partial^r_{n/|n|}F\|_1$. The second form follows from $\partial^r_v=\sum_{|\alpha|=r}\frac{r!}{\alpha!}v^\alpha\partial^\alpha$ and $\sum_{|\alpha|=r}\frac{r!}{\alpha!}|v^\alpha|\le d_0^{r/2}$ for $|v|=1$, and $\sum_{n\ne0}|n|^{-r}<\infty$ for $r>d_0$.
\end{proof}

Nothing in the lemma refers to $\eps$ except through the period and the norms: $C_{d_0,r}$ is absolute, and every $\eps$-dependence sits in the $L^1$ norms of the derivatives. At fixed $\eps$ those norms are finite numbers, continuous in every mass parameter of $F$.

\begin{proof}[Proof of \cref{lem:s10:V}, with the boundedness clause]
Fix $k$. We show that each quantity converges as $J\to\infty$, uniformly for $M\in[M_1,m]$, to its value on $\eps\Z^3$; the limits named in the lemma are these values, and boundedness in $J\ge J_*(k)$ follows from convergence.

\emph{The counterterm.} By (i) of \cref{ssec:a06:setting}, ${\dm}_{(ab),T}(M)$ is defined order by order as minus the sum of the one-particle-irreducible scalar two-point graphs at zero momentum of order $(a,b)$, in which the lower-order counterterms appear as vertex coefficients. Each graph with $L_0$ loops is $\RT{F_G}$ for a smooth $(2\pi/\eps)$-periodic $F_G$ (products of periodic propagators $(\hat p^2+M^2)^{-1}$, $(\hat q^2+\lmu^2)^{-1}$ at the masses $M$, $\lmu$, all positive, times trigonometric polynomials; at fixed $\eps$ there is no singularity); by \cref{lem:a06:RS} with any $r\ge3L_0+1$, $\RT{F_G}\to\int_{\BZ^{L_0}}F_G$ as $\ell\to\infty$ at the rate $\ell^{-r}$, uniformly on $[M_1,m]$ because the $L^1$ norms are continuous in $M$ on that compact interval. By induction on the order, the lower-order coefficients converge, so ${\dm}_{(ab),T}(M)\to{\dm}_{(ab)}(M)$, the coefficient defined by the same recursion with $\int_{\BZ}$ in place of $\RT{\cdot}$, which is King's free-boundary counterterm on $\eps\Z^3$ \cite[p.~330]{K2}: $\alpha^{\mathrm{tor}}_{k,T_J}(M)\to\alpha^\infty_k(M)$ uniformly on $[M_1,m]$. Hence $s_{k,T_J}(M_1)\to s^\infty_k(M_1)$, and the first offset $\sigma_{k,T_J}=\frac12(\alpha^{\mathrm{tor}}_{k,T_J}(m)-\alpha^F)$ tends to $0$, since $\alpha^F=m^2+{\dm}_\eps=\alpha^\infty_k(m)$ (\cref{def:s10:torus}).

\emph{The Wick constant.} $c_{k,T}(M)=\RT{(\hat p^2+M^2)^{-1}}\to\int_{\BZ}(\hat p^2+M^2)^{-1}=\sW_\eps(M)$, and $\Delta c_{T_J}=\sW_\eps-c_{k,T_J}(m)\to0$; this is also the image-sum statement, $c_{k,T}(M)=\sum_{n\in\Z^3}\sW_\eps(M;n\ell)$, each image decaying exponentially at fixed $\eps$.

\emph{The constants.} $E_0(M)/|T|$ is a Riemann sum of the smooth periodic function $\frac N2\ln(\hat p^2+M^2)$ plus an $M$-independent photon term, so $[E_0(M_1)-E_0(m)]/|T|$ converges to $\frac N2\int_{\BZ}\ln\frac{\hat p^2+M_1^2}{\hat p^2+m^2}$. $E_1(M)$ is a finite sum of connected vacuum graphs on $T$; by translation invariance each equals $|T|$ times a Riemann sum $\RT{F_G}$ of a smooth periodic integrand (the graph with one vertex fixed), with the counterterm vertices ${\dm}_{(ab),T}(M)$ as coefficients, so $E_1(M)/|T|$ converges by \cref{lem:a06:RS} and the first step. Hence
\[
|T_J|^{-1}B^z_{k,T_J}(M_1)=\frac{E(M_1)-E(m)}{|T_J|}-\tfrac12\big(2c_{k,T_J}(m)+z_k\big)\big(\alpha^{\mathrm{tor}}_{k,T_J}(m)-\alpha^{\mathrm{tor}}_{k,T_J}(M_1)\big)
\]
converges to a finite limit $b^\infty_k(M_1)$, every term converging and $z_k$ being independent of $J$.
\end{proof}

\subsection{The $c$-number lemma}\label{ssec:a06:T2}

\Cref{hyp:s03:T2} says that $\sigma_1(M_1)=\liminf_ks^\infty_k(M_1)$ and $\beta_1(M_1)=\limsup_kb^\infty_k(M_1)$ are finite. Both are statements about the $\eps\Z^3$ values of \cref{ssec:a06:V}, uniformly in $\eps=\eps_k$, and both come from one source: a mass derivative of King's renormalised two-point graphs lowers every lattice degree by two and leaves a convergent lattice integral.

\begin{proposition}[the derivative identity]\label{prop:a06:D}
Fix $\eps$ and a King torus $T$ (or $\eps\Z^3$, with $\int_{\BZ}$ in place of $\RT{\cdot}$ and densities in place of extensive quantities). Coefficientwise in $(e,\lambda)$,
\[
\frac{d}{dM^2}\,B^z_{k,T}(M)=\tfrac12|T|\,\Psi_{k,T}(M),\qquad
\Psi_{k,T}(M)\defeq\big(1+{\dm}_{k,T}'(M)\big)\big[N\big(c_{k,T}(m)-c_{k,T}(M)\big)+z_k-z_{k,T}(M)\big]-\big[\big(1+{\dm}'_{k,T}(M)\big)\hat Y_{T}(M)\big]_{\le13},
\]
where $'=d/dM^2$, $Y_T(M)\defeq|T|^{-1}\langle\sum_x\eps^3|\phi_x|^2\rangle^{(M)}_{k,T}-Nc_{k,T}(M)$ is the connected one-point function of $|\phi|^2$ in the torus model at mass $M$, $z_{k,T}(M)$ its order-$e^2$ part, $\hat Y_T\defeq Y_T-z_{k,T}(M)$, and $[\cdot]_{\le13}$ the truncation at total order $13$ in $(e,\lambda)$. Consequently $|T|^{-1}|B^z_{k,T}(M_1)|\le\frac12|m^2-M_1^2|\sup_{M\in[M_1,m]}|\Psi_{k,T}(M)|$.
\end{proposition}

\begin{proof}
Write $\ln Z^{(M)}=\mathcal L(M)-E_1(M)$ as in \cref{ssec:a06:setting}. The vertices depend on $M$ only through ${\dm}_{k,T}(M)$, and for the normalised Gaussian measure $\frac{d}{dM^2}\int d\mu_MG=-\frac12\langle G;\sum_x\eps^3|\phi_x|^2\rangle_{\mu_M}$; so $\mathcal L'=-\frac12|T|\big[Y_T+{\dm}'_{k,T}(Nc_{k,T}+Y_T)\big]$. Further $E_0'=-\frac N2|T|c_{k,T}(M)$, and $E_1'=[\mathcal L']_{\le13}$ since $E_1$ is the Taylor part of $\mathcal L$ and the derivatives commute at fixed $\eps$. With $B^z=E(M)-E(m)-\frac12(Nc_{k,T}(m)+z_k)(\alpha(m)-\alpha(M))|T|$ and $\alpha'=1+{\dm}'$,
\[
\frac{dB^z}{dM^2}=\tfrac12|T|\Big\{(1+{\dm}')\big(Nc_{k,T}(m)+z_k\big)-Nc_{k,T}(M)-\big[Y_T+{\dm}'(Nc_{k,T}(M)+Y_T)\big]_{\le13}\Big\}.
\]
The terms ${\dm}'\,Nc_{k,T}(M)$ and $(1+{\dm}')z_{k,T}(M)$ have order at most $6$ and pass through the truncation unchanged (${\dm}$ has orders $a+2b\le4$); collecting gives $\Psi$. Odd powers of $e$ vanish by charge conjugation, and the order-$\lambda$ part of $Y_T$ vanishes because ${\dm}_{(01),T}=-2\lambda(N+2)c_{k,T}(M)$ uses the same torus tadpole. Since $B^z(m)=0$, $B^z(M_1)=-\frac12|T|\int_{M_1^2}^{m^2}\Psi\,d\mu^2$, which gives the last claim.
\end{proof}

By \cref{lem:s10:V}, $b^\infty_k(M_1)=\lim_J|T_J|^{-1}B^z_{k,T_J}(M_1)=-\frac12\int_{M_1^2}^{m^2}\Psi^\infty_k(\mu)\,d\mu^2$ with $\Psi^\infty_k$ the $\eps\Z^3$ form of $\Psi$, so $\beta_1(M_1)<\infty$ follows from a bound on $\Psi^\infty_k$ uniform in $k$ and in $M\in[M_1,m]$; and $s^\infty_k(M_1)=\frac12(m^2-M_1^2)+\frac12\int_{M_1^2}^{m^2}{\dm}_k'(\mu)\,d\mu^2$, so $\sigma_1(M_1)$ is controlled by the mass derivative of the counterterm. Both are finite sums of renormalised lattice Feynman integrals at zero external momentum, and the tool is Reisz's power counting.

\paragraph{Reisz's theorems in three dimensions.} Reisz's power counting theorem \cite[Thm.~1]{Rei88} states that a lattice Feynman integrand of his class $\mathcal F$ (numerators trigonometric polynomials in the momenta, denominators $e(l;\eps)+\mu^2$ with $e(l;\eps)=\eps^{-2}\eta(l\eps)$, $\eta\ge0$, $\eta(x)\to x^2$) whose lattice degrees of divergence (ultraviolet degrees in the sense of his Definition~2.3, computed with the lattice behaviour of the propagators) are negative for the integral and for every Zimmermann subspace has a convergent continuum limit, and its proof gives the absolute bound $\sup_{\eps\le\eps_0}\int_{\BZ^{L_0}}|F|\le K$. His renormalisation theorem \cite{Rei88b} states that the forest formula with his periodic subtraction operators, applied to a graph all of whose one-particle-irreducible subgraphs have degree $\omega\ge0$ subtracted to order $\omega$, produces an integrand in $\mathcal F$, periodic in all loop momenta, with all Zimmermann degrees negative, provided the degrees are additive over disjoint subgraphs and the subtraction degrees satisfy his conditions (4-3) and (4-4) \cite[Lemma~5.2]{Rei88b}. Both are printed for $d=4$; the dimension enters only through the degrees, and the proofs apply verbatim in $d=3$ with the degrees $\omega=3L_0-2I_\phi-2I_A+(\text{numerator degree})$. King's graphs satisfy the hypotheses: his vertices are the quartic vertex, the hopping vertices with one or two photon legs (trigonometric numerators), and the counterterm vertices; the only divergent one-particle-irreducible subgraphs are the scalar two-point function ($\omega=1$ at one loop, $\omega=0$ at two loops) and the photon two-point function ($\omega=1$), and King's recursive counterterm is exactly the forest-scheme counterterm: on $\eps\Z^3$ the order-$1$ Taylor term of the scalar two-point subtraction vanishes by parity and the photon two-point function at zero momentum vanishes by the lattice Ward identity, so the subtracted integrands are King's graphs with his ${\dm}_{(ab)}$ inserted. We write $\hat RF_G$ for the subtracted integrand of a graph $G$ with these rules, and $\hat R'F_{G,l}$ for the mass-derivative integrand defined next.

\begin{lemma}[derivative integrands]\label{lem:a06:Dp}
Let $G$ be one of King's two-point or vacuum graphs on $\eps\Z^3$ with its scalar two-point subgraphs subtracted to order $1$ (one loop) or $0$ (two loops) and its photon two-point subgraphs to order $1$. \textup{(i)} $\frac{d}{dM^2}\int_{\BZ^{L_0}}\hat RF_G=\sum_{l}\int_{\BZ^{L_0}}\hat R'F_{G,l}$, the sum over the scalar lines $l$ of $G$, where $F_{G,l}$ has the line $l$ doubled (a factor $-(\hat p_l^2+M^2)^{-1}$) and $\hat R'$ uses the same forests with the \emph{effective} degrees: unchanged for subgraphs not containing $l$; $0$ (the order-$1$ Taylor term dropped) for a one-loop two-point subgraph containing $l$; no subtraction for a two-loop two-point subgraph containing $l$. \textup{(ii)} $\hat R'F_{G,l}$ lies in $\mathcal F$, is periodic in all loop momenta, and has all Zimmermann degrees negative; the same holds for every momentum derivative $\partial^\alpha_p\hat R'F_{G,l}$ and $\partial^\alpha_p\hat RF_G$, with no degree increased.
\end{lemma}

\begin{proof}
(i) The subtraction operators act on the external momentum of a subgraph and do not depend on $M$, so $d/dM^2$ commutes with them and acts on the propagators, $\frac{d}{dM^2}(\hat p^2+M^2)^{-1}=-(\hat p^2+M^2)^{-2}$: this gives $\sum_l\hat R^{\mathrm{formal}}F_{G,l}$ with the original degrees. Doubling a line lowers the degree of every subgraph containing it by $2$, so a one-loop two-point subgraph containing $l$ has degree $-1$ and a two-loop one degree $-2$; the Taylor terms of order $>$ (new degree) are over-subtractions, which Reisz's conditions do not allow. They integrate to zero: the order-$1$ term of a one-loop subgraph is $\sum_\rho P(p_\rho;\eps)\,\partial_{p_\rho}F_\gamma(q,0)$ with $q$ the loop momentum of the subgraph, a free integration variable at fixed outer momenta, and $\partial_pF_\gamma(q,0)$ is odd under $q\to-q$ while the rest of the integrand does not depend on $q$; the order-$0$ term of a two-loop subgraph containing the doubled line is its value at zero momentum, which is the mass derivative of the two-loop counterterm and is cancelled exactly by the derivative of the counterterm vertex that the forest formula carries. Dropping them gives $\hat R'$. (ii) With the effective degrees, every subtracted subgraph is subtracted to its (new) degree and every unsubtracted one has negative degree, which is Reisz's condition (4-3) with equality and (4-4); his Lemma 5.2 gives negative Zimmermann degrees. Momentum derivatives lower the degree of the differentiated propagator by one and keep the class $\mathcal F$ and the periodicity.
\end{proof}

\begin{proposition}[uniform bounds]\label{prop:a06:RD}
For each of King's graphs $G$, each scalar line $l$ and each multi-index $\alpha$ there are $\eps_0>0$ and $K<\infty$ with
\[
\sup_{\eps\le\eps_0}\ \sup_{M\in[M_1,m]}\ \int_{\BZ^{L_0}}\big|\partial^\alpha_p\hat RF_G(p;M)\big|\,dp\le K,\qquad
\sup_{\eps\le\eps_0}\ \sup_{M\in[M_1,m]}\ \int_{\BZ^{L_0}}\big|\partial^\alpha_p\hat R'F_{G,l}(p;M)\big|\,dp\le K,
\]
and for $\eps_0<\eps\le1/5$ the same integrals are bounded with an $\eps_0$-dependent constant.
\end{proposition}

\begin{proof}
By \cref{lem:a06:Dp} the integrands are in $\mathcal F$ with negative Zimmermann degrees, so Reisz's theorem gives the absolute bound at each fixed mass. Uniformity on the compact mass interval: every term of the integrand is a mass-free trigonometric numerator over a product of propagators $(e_L+\mu_L^2)^{-n_L}$; bringing the terms to the common denominator $\prod_L(e_L+\mu_L^2)^{n_L}$, the numerator is a polynomial in the squared masses with coefficients that are mass-free trigonometric polynomials, $\sum_iP_i(\mu^2)V_i(p;\eps)$ with finitely many distinct monomials $P_i$; each $V_i/\prod_L(e_L+\mu_L^2)^{n_L}$ has the same or lower degrees, so is bounded in $L^1$ by Reisz's theorem at the fixed mass $\mu_{\min}=\min(M_1,\lmu)$, and dominates the integrand at every mass in the interval up to the factor $\sup_i\sup_{[M_1,m]}|P_i|\prod_L(\mu_{\min}^2/(\ldots))$; since all masses are bounded above and below on the interval, the constant is uniform. For $\eps>\eps_0$ the integrals are finite numbers of a finite family, continuous in $M$.
\end{proof}

\begin{theorem}[the $c$-number lemma]\label{thm:a06:CD1}
Let $0<M_1<m$. \textup{(a)} There is $C=C(e,\lambda,\lmu,m,M_1)$ with $\sup_{k}\sup_{M\in[M_1,m]}|\Psi^\infty_k(M)|\le C$; hence $\sup_k|b^\infty_k(M_1)|\le\frac12(m^2-M_1^2)C$ and $\beta_1(M_1)<\infty$. \textup{(b)} For every $M>0$ the limit $\alpha_\infty(M)-\alpha_\infty(m)=\lim_k[\alpha^\infty_k(M)-\alpha^\infty_k(m)]$ exists and is a polynomial of degree $\le2$ in $(e^2,\lambda)$ with finite coefficients; hence $\sigma_1(M_1)=\lim_ks^\infty_k(M_1)=\frac12(\alpha_\infty(m)-\alpha_\infty(M_1))$ exists and is finite. In particular \cref{hyp:s03:T2} holds, and so does the statement \sxivCT\ of the coupling-range section (before \cref{lem:s14:CW}).
\end{theorem}

\begin{proof}
(a) By \cref{prop:a06:D}, $\Psi^\infty_k$ is a polynomial of degree $\le13$ in $(e,\lambda)$ whose coefficients are finite combinations of four kinds of quantities, each to be bounded uniformly in $\eps=\eps_k\le\eps_0$ and $M\in[M_1,m]$: (Q1) the coefficients $x_{ab}(M)$ of $\hat Y(M)$ at orders $a/2+b\ge2$; (Q2) the coefficients $y_{a'b'}(M)$ of ${\dm}_k'(M)$ at the five orders $(0,1),(2,0),(0,2),(2,1),(4,0)$; (Q3) $c(m)-c(M)=\sW_\eps(m)-\sW_\eps(M)$; (Q4) $z_k-z_k(M)$. (Q1): $x_{ab}(M)$ is a sum over the connected two-point graphs of $\langle|\phi|^2\rangle$ (rings of scalar lines with the vertices of the action, the subdivergences subtracted by King's counterterms) of $\int_{\BZ^{L_0}}\hat RF_G$ with overall degree $\omega=3-2(a/2+b)-\ldots\le-1$ and negative proper degrees; \cref{prop:a06:RD} with $\alpha=0$ bounds each. (Q2): $y_{a'b'}(M)=-\frac{d}{dM^2}$ of the zero-momentum one-particle-irreducible two-point function at order $(a',b')$ without its own counterterm $=-\sum_l\int\hat R'F_{G,l}$ by \cref{lem:a06:Dp}(i), bounded by \cref{prop:a06:RD}. (Q3): $\sW_\eps(m)-\sW_\eps(M)=(M^2-m^2)\int_{\BZ}(\hat p^2+m^2)^{-1}(\hat p^2+M^2)^{-1}$, a degree $-1$ integral, bounded. (Q4): $z_k=z_k(m)$ is the reference scheme of \cref{rem:s02:scheme}, the order-$e^2$ coefficient of $Y$ at mass $m$ (\cref{lem:k02:cnum}); $z_k(m)-z_k(M)=-\int_{M^2}^{m^2}\frac{d}{d\mu^2}z_k(\mu)\,d\mu^2$, and $\frac{d}{d\mu^2}z_k(\mu)=e^2\sum_l\int\hat R'F_{(20),l}$ is bounded by \cref{prop:a06:RD} (the degree-$0$ one-loop photon-exchange and seagull rings, each lowered to $-2$ by the doubled line); integrate over the interval. For the finitely many $k$ with $\eps_k>\eps_0$ each quantity is a finite number, continuous in $M$. (b) For each order $(a',b')$, $y_{a'b'}(\mu;\eps)=-\sum_l\int_{\BZ^{L_0}}\hat R'F_{G,l}$ converges as $\eps\to0$ by Reisz's renormalisation theorem applied to the effective integrand (\cref{lem:a06:Dp}(ii) provides its hypotheses) and is bounded uniformly in $\mu$ on compact intervals (\cref{prop:a06:RD}); dominated convergence in ${\dm}_k(m)-{\dm}_k(M)=\int_{M^2}^{m^2}{\dm}_k'(\mu)\,d\mu^2$ gives the limit, a polynomial of degree $\le2$ in $(e^2,\lambda)$ since ${\dm}$ has orders $a+2b\le4$ and odd powers of $e$ vanish. The divergent parts of the counterterm, $1/\eps$ at one loop and $\ln(1/\eps)$ at two loops, are mass-independent and cancel in the difference; this is what the negative degrees of the derivative integrands express.
\end{proof}

The one-loop evaluation of the limit in (b) is the closed form of \cref{lem:s14:CW}(b)(ii); the two-loop coefficients are the finite numbers of (b), bounded in \cref{prop:s14:T3} by explicit majorants. Nothing in this subsection uses the torus beyond \cref{lem:s10:V}: the torus forms of (Q1)--(Q4), with $\RT{\cdot}$ in place of $\int_{\BZ}$ and the additional holonomy vertices (the torus photon two-point function at zero momentum, which vanishes on $\eps\Z^3$ by the Ward identity but is $O(\ell^{-r})$ on $T$ by \cref{lem:a06:RS}), are bounded uniformly in $\eps$ and in the torus by the same argument, but the Wick leg does not need them.
